% Part: first-order-logic
% Chapter: tableaux
% Section: soundness-identity

\documentclass[../../../include/open-logic-section]{subfiles}

\begin{document}

\olfileid{fol}{tab}{sid}

\olsection{السلامة مع \usetoken{S}{identity}}

\begin{prop}
!!^{tableau}s ذات قواعد الهوية سليمة: فلا !!{tableau} مغلق قابل للإرضاء.
\end{prop}

\begin{proof}
يكفي، على ما تقدّم، أن نأخذ فرعًا قابلًا للإرضاء من !!a{tableau}
ونثبت أن إجراء إحدى قواعد~$\eq$ عليه يبقيه قابلًا للإرضاء.
نسمّي~${\OLId{greek-variable}{Gamma}}$ مجموعة ال!!{formula}s الموقّعة على الفرع،
ونأخذ~$\Struct{M}$، وهي !!a{structure} ترضي~${\OLId{greek-variable}{Gamma}}$.

أما إن مدّدنا الفرع بقاعدة~$\eq$، فقد أضفنا ال!!{formula}
الموقّعة~$\sFmla{\True}{\eq[{\OLId{latin-ordinary}{t}}][{\OLId{latin-ordinary}{t}}]}$. وصدق
$\Sat{{\OLId{latin-looped}{M}}}{\eq[{\OLId{latin-ordinary}{t}}][{\OLId{latin-ordinary}{t}}]}$ بيّن؛ فتبقى~$\Struct{M}$ مرضيةً
لـ~${\OLId{greek-variable}{Gamma}} \cup
\{\sFmla{\True}{\eq[{\OLId{latin-ordinary}{t}}][{\OLId{latin-ordinary}{t}}]}\}$.

وأما إن مدّدناه بـ~$\TRule{\True}{\eq}$، فنضيف !!{formula}
موقّعة هي~$\sFmla{\True}{!A({\OLId{latin-ordinary}{t}}_{\OLMathNumeral{2}})}$، وقد اشتملت~${\OLId{greek-variable}{Gamma}}$ على
$\sFmla{\True}{\eq[{\OLId{latin-ordinary}{t}}_{\OLMathNumeral{1}}][{\OLId{latin-ordinary}{t}}_{\OLMathNumeral{2}}]}$ و$\sFmla{\True}{!A({\OLId{latin-ordinary}{t}}_{\OLMathNumeral{1}})}$.
فيلزم~$\Sat{{\OLId{latin-looped}{M}}}{\eq[{\OLId{latin-ordinary}{t}}_{\OLMathNumeral{1}}][{\OLId{latin-ordinary}{t}}_{\OLMathNumeral{2}}]}$ و$\Sat{{\OLId{latin-looped}{M}}}{!A({\OLId{latin-ordinary}{t}}_{\OLMathNumeral{1}})}$.
نأخذ~${\OLId{latin-ordinary}{s}}$ إسنادًا يحقق~${\OLId{latin-ordinary}{s}}({\OLId{latin-ordinary}{x}}) = \Value{{\OLId{latin-ordinary}{t}}_{\OLMathNumeral{1}}}{{\OLId{latin-looped}{M}}}$.
وبمقتضى \olref[syn][ass]{prop:sentence-sat-true}
نحصل على~$\Sat{{\OLId{latin-looped}{M}}}{!A({\OLId{latin-ordinary}{t}}_{\OLMathNumeral{1}})}[{\OLId{latin-ordinary}{s}}]$. ولأن~$\varAssign{{\OLId{latin-ordinary}{s}}}{{\OLId{latin-ordinary}{s}}}{{\OLId{latin-ordinary}{x}}}$،
تفيد \olref[syn][ext]{prop:ext-formulas} أن~$\Sat{{\OLId{latin-looped}{M}}}{!A({\OLId{latin-ordinary}{x}})}[{\OLId{latin-ordinary}{s}}]$.
ومن~$\Sat{{\OLId{latin-looped}{M}}}{\eq[{\OLId{latin-ordinary}{t}}_{\OLMathNumeral{1}}][{\OLId{latin-ordinary}{t}}_{\OLMathNumeral{2}}]}$ يلزم~$\Value{{\OLId{latin-ordinary}{t}}_{\OLMathNumeral{1}}}{{\OLId{latin-looped}{M}}} = \Value{{\OLId{latin-ordinary}{t}}_{\OLMathNumeral{2}}}{{\OLId{latin-looped}{M}}}$،
فيكون~${\OLId{latin-ordinary}{s}}({\OLId{latin-ordinary}{x}}) = \Value{{\OLId{latin-ordinary}{t}}_{\OLMathNumeral{2}}}{{\OLId{latin-looped}{M}}}$. نعيد استعمال
\olref[syn][ext]{prop:ext-formulas}، فيحصل
$\Sat{{\OLId{latin-looped}{M}}}{!A({\OLId{latin-ordinary}{t}}_{\OLMathNumeral{2}})}[{\OLId{latin-ordinary}{s}}]$. ثم ننتهي، بمقتضى
\olref[syn][ass]{prop:sentence-sat-true}، إلى~$\Sat{{\OLId{latin-looped}{M}}}{!A({\OLId{latin-ordinary}{t}}_{\OLMathNumeral{2}})}$.
وحالة~$\TRule{\False}{\eq}$ تجري على هذا المسلك نفسه.
\end{proof}

\end{document}
